\section{Fast-Forwarding Dykstra's Algorithm}
\label{sec:main_result}

% In this section, we propose a modification to Dykstra's algorithm to resolve the stalling problem in the general case.

% \subsection{Stalling Formalisation}
Based on the discussion of stalling in~\cite{DYKSTRASTALLING}, we identify as a necessary condition for stalling of the iterates for successive iteration cycles over $n$ sets. This is formalised in Definition~\ref{def:stalling}.

\begin{defn}[Stalling]\label{def:stalling}
Given $m\geq n-1$, a stalling period is defined as those $i\in\lbrace 0,\dots,N_{stall}\rbrace$ with $N_{stall}\geq n$ for which $x_{m+i}=x_{m+i-n}$.
\end{defn}

During stalling, only the auxiliary variables $e_{m+i}$ change, which are modified by constant increments $x_m-x_{m+1}$. The stalling period therefore continues until $e_{m+i}$ becomes such that $x_{m+i}+e_{m-n+i}\in\mathcal{H}_{[m+i]}$. For polyhedral sets, the length of the stalling period can be pre-computed as soon as the algorithm encounters stalling, using the half-spaces that are active at this iteration. The stalling period continues until one of these half-spaces becomes inactive, at which point stalling terminates. This is formalised in Theorem~\ref{thm:nstall}.

\begin{thm}[Length of Stalling Period]\label{thm:nstall}
Suppose that stalling starts at iteration $m$ as in Definition~\ref{def:stalling} and denote the active half-spaces by $\mathcal{A}\coloneqq\set{i\in\lbrace 0,\dots,n-1\rbrace}{x_{m+i}+e_{m-n+i}\not\in\mathcal{H}_{[m+i]}}$. Then, the length of the stalling period is given by
\begin{align}\label{eq:nstall}
N_{stall}\coloneqq\min_{i\in\mathcal{S}} \left(i+\left\lceil k_{m-n+i} / \left(n(c_{[m+i]}-x_{m+i}^T f_{[m+i]})\right)\right\rceil\right),
\end{align}
where $k_{m-n+i}=e_{m-n+i}^T f_{[m+i]}$ and $\mathcal{S}\coloneqq\set{i\in \mathcal{A}}{x_{m-n+i}^T f_{[m+i]} - c_{[m+i]} < 0}$.
\end{thm}

\begin{pf}
Note that the existence of a finite number $N_{stall}$ is a consequence of the Boyle-Dykstra theorem. According to~\eqref{eq:dykstra:proj:poly}, the stalling period will terminate once an active half-space becomes inactive, i.e.\ $x_{m+i}+e_{m-n+i}\in\mathcal{H}_{[m+i]}$ for some $i\geq n$ (Definition~\ref{def:stalling}). According to~\eqref{eq:km}, the only half-spaces that can become inactive are those for which $x_{m+i}^T f_{[m+i]} - c_{[m+i]} < 0$, which is, by Definition~\ref{def:stalling}, a constant during stalling. The length of the stalling period can therefore be obtained from choosing the smallest possible integer $N_i$ for which
\begin{align*}
k_{m-n+i}+n N_i\left(x_{m+i}^T f_{[m+i]} - c_{[m+i]}\right) < 0,\qquad i\in\mathcal{A},
\end{align*}
so that $N_i=\left\lceil k_{m-n+i} / \left(n(c_{[m+i]}-x_{m+i}^T f_{[m+i]})\right)\right\rceil$. If $i_1<\dots<i_j$ are such that $N_{i_1}=\dots = N_{i_j}$, then according to~\eqref{eq:dykstra}, half-space $i_1$ is discarded first (and thus the minimiser of~\eqref{eq:nstall}).\qed
\end{pf}
Following the end of the stalling period, one half-space is discarded from the active half-space and $\mathcal{A}$ redefined, upon which the algorithm may again encounter a stalling condition. By Theorem~\eqref{thm:nstall}, the stalling period can be fast-forwarded by applying constant increments to $k_{m+i}$ for $i=0,\dots,N_{stall}$.

Theorem~\ref{thm:nstall} and Dykstra's algorithm can be combined as in Algorithm~\ref{alg:fastforward}. The algorithm proceeds identically to the standard Dykstra's algorithm until it detects stalling by verifying the stationarity condition $x_m = x_{m-n}$, as per Definition~\ref{def:stalling}. Once stalling is detected at iteration $m$, the algorithm executes its fast-forwarding phase. After performing the standard variable update for the current iteration, it computes the exact number of required stalling iterations, $N_{stall}$, by calling a function that implements Theorem~\ref{thm:nstall}. From this, it determines the number of full $n$-step cycles to skip, $N_{\text{cycles}} \coloneqq \lceil N_{stall} / n \rceil$. For each of the $n$ dual variables $e_j$, the algorithm then computes the constant increment $\Delta k_j = (\trans{f_j} x_{\text{stall}} - c_j)$ and applies a single, large update: $e_j \leftarrow e_j + N_{\text{cycles}} \cdot (\Delta k_j f_j)$. This ``jumps'' the dual variables forward to their state at iteration $m + N_{stall}$. Finally, the main iteration counter $m$ is advanced by $N_{stall}$, effectively skipping the entire stalling period. 

Since the proposed algorithm achieves the same iterates after escaping the stalling period (lines 13-18) as would be attained with standard Dykstra, the Boyle-Dykstra theorem~\cite{DYKSTRA} is retained, thus the iterates are guaranteed to converge to the optimal solution $x^\star$ asymptotically. This is formalised in Corollary~\ref{cor:convergence} below.

\begin{cor}[Convergence of Algorithm~\ref{alg:fastforward}]
\label{cor:convergence}
The sequence of iterates $\{x_m\}$ generated by the Stall-Averse Dykstra's Algorithm (Algorithm~\ref{alg:fastforward}) converges asymptotically to the optimal solution $x^\star = \mathcal{P}_{\mathcal{H}}(x^\circ)$, thus
\begin{align} \label{eq:boyle-dykstra}
    \lim_{m\rightarrow\infty}\|x_m-x^\star\|_2=0
\end{align}
\end{cor}

\begin{algorithm}
\caption{Stall-Averse Dykstra's Algorithm}
\label{alg:fastforward}
\begin{algorithmic}[1]
    \Function{DykstraFastForward}{$x^\circ, \{\mathcal{H}_i\}_{i=0}^{n-1}, N_{\max}, \epsilon_{\text{stall}}$}
        \State $x \leftarrow x^\circ$; $e_j \leftarrow 0$ \textbf{for} $j=0, \dots, n-1$; $X_{\text{hist}}[0] \leftarrow x^\circ$; $m \leftarrow 0$ \algcomment{Initialisation}
        \Statex\hrulefill
        \While{$m < N_{\max}$}
            \algcomment{Standard Dykstra's update}
            \State $e_{[m]} \leftarrow x + e_{[m-n]} - \mathcal{P}_{[m]}(x + e_{[m-n]})$
            \State $x \leftarrow x + e_{[m-n]} - e_{[m]}$ 
            \State $X_{\text{hist}}[m+1] \leftarrow x$
            \Statex\hrulefill
            \If{$m \ge n \wedge \twonorm{x - X_{\text{hist}}[m-n]} < \epsilon_{\text{stall}}$} \algcomment{Stalling detected per Definition~\ref{def:stalling}}
                \State $x_{\text{stall}} \leftarrow x$
                \State $N_{stall} \leftarrow \min_{j \in \mathcal{S}} \left(j+\left\lceil k_j / \left(n(c_j - \trans{f_j} x_{\text{stall}})\right)\right\rceil\right)$ \algcomment{Compute using Thm.~\ref{thm:nstall}}
                \State $N_{\text{cycles}} \leftarrow \lceil N_{stall} / n \rceil$
                \For{$j \in \{0, \dots, n-1\}$}
                    \State $k_j \leftarrow \trans{f_j} e_j$ \algcomment{Get scalar from $e_j = k_j f_j$}
                    \State $\Delta k_j \leftarrow \trans{f_j} x_{\text{stall}} - c_j$
                    \State $e_j \leftarrow (k_j + N_{\text{cycles}} \times \Delta k_j) f_j$ \algcomment{Fast-forward dual variable}
                \EndFor
                \State $m \leftarrow m + N_{stall}$ \algcomment{Jump iteration counter}
            \Else
                \State $m \leftarrow m + 1$ \algcomment{Increment iteration counter}
            \EndIf
            \Statex\hrulefill
        \EndWhile
        \State \Return $x$
    \EndFunction
\end{algorithmic}
\end{algorithm}

% \begin{cor}[Convergence of Algorithm~\ref{alg:fastforward}]
% \label{cor:convergence}
% The sequence of iterates $\{x_m\}$ generated by the Stall-Averse Dykstra's Algorithm (Algorithm~\ref{alg:fastforward}) is guaranteed to converge to the optimal solution $x^\star$ asymptotically.
% \end{cor}
% \begin{pf}
% The fast-forwarding step (lines 13--18) is not an approximation. It computes the exact state of the dual variables $\{e_j\}$ that would be attained by the standard Dykstra's algorithm~\eqref{eq:dykstra} after $N_{stall}$ iterations. Since the proposed algorithm achieves the same iterates after escaping the stalling period as would be attained with standard Dykstra, it retains the convergence guarantees of the original Boyle-Dykstra theorem~\cite{DYKSTRA,DYKSTRAPERKINS}. \qed
% \end{pf}